\documentclass[11pt]{article}
\usepackage[a4paper,margin=2.5cm]{geometry}
\usepackage[T1]{fontenc}
\usepackage{lmodern}
\usepackage[hidelinks]{hyperref}
\setlength{\parindent}{0pt}
\setlength{\parskip}{0.7em}
\pagestyle{empty}

\begin{document}

\begin{flushright}
27 October 2026
\end{flushright}

To the Editors, \emph{Physics of the Dark Universe}

Dear Editors,

We submit our manuscript \emph{``The Bondi dipole in full numerical relativity: a self-accelerating positive--negative mass binary''}, by Nikita M.\ Shirokov and Ilya Nachevsky, for consideration as a research article in \emph{Physics of the Dark Universe}.

Negative mass recurs in the physics of the dark sector, and this journal has recently published three studies of it that our manuscript cites: whether negative-mass objects can exist in the sky (Nojiri and Odintsov, PDU \textbf{53}, 102364, 2026), the gravitational-wave signals of negative-mass binaries (Trivedi and Loeb, PDU \textbf{53}, 102409, 2026), and whether negative mass can drive the cosmological expansion (Manfredi, Rouet and Miller, PDU 102493, 2026). Bondi showed in 1957 that a body of negative and a body of positive mass should accelerate together without limit while their total momentum stays zero. Our manuscript tests this prediction, to our knowledge for the first time, in full nonlinear general relativity. Its main results are:

\begin{itemize}\setlength{\itemsep}{0.2em}
\item We construct constraint-satisfying initial data for stars of a complex scalar field that enters the Einstein equations with the opposite sign. These stars have positive inertial and passive masses but negative active and ADM masses, and are, to our knowledge, the first asymptotically flat bodies of negative ADM mass evolved in numerical relativity.
\item A pair of stars with equal and opposite masses, released from rest, accelerates as a single object at constant separation, and its total momentum stays zero to within 0.5\%.
\item The acceleration follows Newton's law of gravitation in both mass and separation, reverses when the two stars are exchanged, and is absent in same-sign pairs and in isolated stars. A resolution study shows the constraints and the measured motion converging.
\item We detect no gravitational radiation. In the longest run the pair reaches a tenth of the speed of light with both stars intact.
\end{itemize}

The work computes from the full theory the dynamics that such studies take as input, and it identifies what is needed to reach the relativistic speeds at which the pair would become a source of the signals proposed by Trivedi and Loeb. We therefore believe it suits the journal's readership.

The evolution code, the initial-data code, the parameter files and the time series of every simulation are public (archived at \url{https://doi.org/10.5281/zenodo.23662692}), and every figure and quoted number regenerates from them without a GPU. A preprint is available as arXiv:2608.24577. The manuscript has not been published and is not under consideration elsewhere, and both authors have approved this submission.

Yours sincerely,

Nikita M.\ Shirokov (corresponding author)\\
Moscow Institute of Physics and Technology\\
\href{mailto:shirokov.nm@phystech.edu}{shirokov.nm@phystech.edu}

\end{document}
